% Section 9 -- Related Work
\section{Related Work}
\label{sec:related}

This section traces three lines of work that converge on joint BGP-traffic anomaly attribution: the evolution of routing anomaly detection, the statistical and counterfactual machinery for causal attribution, and the dependability discipline that constrains how such systems must be evaluated.

\subsection{Routing anomaly detection}

Early routing anomaly detection worked entirely in the control plane, comparing announcements against a maintained expectation of who should be announcing a prefix. ARTEMIS made this efficient by signing the routing table and validating at line rate~\cite{artemis}. The limitation was structural: control-plane observation alone cannot distinguish an announcement that degraded reachability from one that was merely unusual, and it emits an unordered candidate stream that still requires operator adjudication.

A second line moved to the data plane. RECON and Netalyzr probe reachability from distributed vantage points and infer reverse paths, answering whether connectivity actually broke~\cite{netalyzr-2013}. Their coverage is bounded by vantage-point count and cost grows with address space, so they complement rather than replace control-plane monitoring.

Hybrid systems then began fusing the two. Pheme and Argus correlate control-plane events with data-plane measurements to suppress false positives~\cite{pheme-2020,argus-2019}, and the BGP-graph community has since produced learned detectors that score prefixes by AS-path and update-stream features~\cite{keegan-2023-bgp-graph,beam2024}. These improve per-event precision but inherit the supervision problem: they need labeled incidents, of which the public record contains few, and their outputs are scores rather than auditable arguments. Stream4BoT addresses the throughput dimension but not auditability~\cite{stream4r-2020}.

A parallel line attacked route leaks specifically. The valley-free violation is the right structural signature of a leak; Peerlock and KRELAY turned that into an enforcement mechanism between consenting ASes~\cite{peerlock,krerelay-2019}. Because participation is voluntary and the mechanism protects a single leak class, these systems reduce incidence but do not provide attribution.

The RPKI line addressed origin validation cryptographically. Relying parties verify signatures against the RPKI profile; measurements quantify what remains unprotected~\cite{cooper-2013-rpki,testart-2025-rpki,ahmad-2024-rpki-fail}. RoST and ImpROV characterize the collateral damage that valid-but-wrong announcements still cause~\cite{rost2025,improv2025}. The common finding is that deployment is partial and incremental, so origin validation reduces but does not remove the need for attribution.

\subsection{Joint correlation and causal attribution}

The statistical line began with Granger, who asked whether the past of one series improves prediction of the present of another~\cite{granger-1969-investigating}. Granger causality is the natural starting point for BGP-traffic joint analysis, and in our measurements it is also the weakest: the test requires stationary residuals, and BGP update streams are bursty and non-stationary by construction, so the precondition is systematically violated.

Transfer entropy relaxed Granger's linearity assumption by working in the information domain and brought it to network telemetry~\cite{schreiber-2000-te,barnett-2009-granger-network}. It improves on Granger but inherits a second problem: the estimator is sensitive to embedding dimension and history length, and it yields a single scalar with no notion of when the evidence is too weak to support a claim.

The counterfactual literature asked a strictly stronger question. Pearl's probability of necessity and sufficiency asks whether treatment was necessary for the effect and sufficient to produce it~\cite{pearl-causality-2009,pearl1999}. The identification results were sharpened by Shpitser and Pearl and extended to graph-free settings, making PNS estimable exactly when a suitable causal graph and adjustment set exist~\cite{shpitser-2008-pns,pearl-2014-pns-3}. The gap for network telemetry is that no a priori causal graph exists: BGP updates and traffic are jointly produced by the same operator decisions, so the graph the theorem requires is exactly what the analyst lacks. Graph-free estimators reduce to difference-in-means on a pre/post partition.

Recent network-operations work has begun importing this machinery. Liu et al. learn a causal graph over BGP telemetry and use it to attribute BGP anomaly chains~\cite{liu-2024-bgp-cause}; Cho et al. apply Bayesian discovery to BGP attribution directly~\cite{cho-2025-causal-bgp}; and Hu et al. apply counterfactual estimation to traffic anomalies~\cite{hu-2025-traffic-causal}. All three confront the same missing-graph problem and settle on a timeline-partition estimator at the routing breakpoint. Their designs differ from ours in a critical respect: they report a point estimate; we require the estimator's significance conditions to hold before the engine will commit, and abstain otherwise. The abstention discipline is what converts the estimator from a score into a falsifiable claim.

\subsection{Dependability and system papers}

The dependability community has framed anomaly attribution as a fault-tolerance problem, with reproducibility as a first-class review criterion. Because evaluation sets here are necessarily small, the Wilson interval, the bootstrap, and permutation tests remain the workhorse inference tools.

System papers have contributed the architectural vocabulary. Multi-agent orchestration in the ReAct style popularized the blackboard pattern in which specialized reasoners share a structured cache, which is the arrangement Section~\ref{sec:overview} adopts. BGP semantics-aware representation learning showed that embedding update streams and operator-visible context jointly predict operator response better than update volume alone. Testbed-grounded inference over forwarding behaviour established that data-plane and control-plane anomalies are frequently independent, which is the empirical basis for treating them as separate evidence channels. Deployment studies of RPKI reliance quantified how much risk survives partial deployment.

What none of these lines provides is a rule for combining control-plane and data-plane evidence into a single adjudicated verdict. BGP Tracer's contribution is that rule: Phase 1 produces a five-class correlation verdict (Section~\ref{sec:overview-phase1}) that tells operators whether a routing anomaly has traffic impact; Phase 2 then performs causal attribution (Section~\ref{sec:overview-phase2}) to determine the cause type and direction. Every prior system has produced either a ranked list of candidates, a single-plane answer, or a non-reproducible operator judgement. BGP Tracer is the first deterministic, reproducible, auditable joint-attribution system for BGP-telemetry anomalies. The dependability discipline enforces reproducibility at every step: the cache is the artifact, the evaluation is the audit, and the abstention record is the interpretability proof~\cite{wilson-1927,efron-1979-bootstrap,phipson-2010-permutation}.
